Orbiter can compute the orbits of a group on 
the lattice of subspaces of a finite vector space.

\bigskip



The orthogonal group is the stabilizer of a non-degenerate quadric.  
Suppose we want to classify the subspaces in $\PG(3,2)$ under 
the action of the orthogonal group. 
In $\PG(3,2),$ there are two distinct nondegenerate quadrics, 
$Q^+(3,2)$ and $Q^-(3,2).$ 
The $Q^+(3,2)$ quadric is a given by the equation 
$$
x_0x_1 + x_2x_3 = 0.
$$
The stabilizer of the quadric is the group $\PGO^+(4,2).$
Suppose we want to classify the orbits of this group on the underlying space $\PG(3,2).$ 
The command

\sourcecode{subspaces_Op_4_2}

produces a classification of all 
subspaces of $\PG(3,2)$ under $\PGO^+(4,2).$ The poset of orbits is shown below.
Each node represents one orbit on subspaces:

\orbiteroutput{GRAPHICS/WRAPPERS/Op_4_3_orbits_drawing}

The nodes show the ranks of points in $\PG(3,2)$ 
as described in Section~\ref{sec:FPG}.
It is also possible to expand the poset to the level of orbit elements. 
The command also produces a detailed view of the poset, showing every subspace. 
Orbits of subspaces are indicated by grouping the corresponding nodes 
at a horizontal level:

\orbiteroutput{GRAPHICS/WRAPPERS/Op_4_2_poset_full_lvl_4_draw}

\bigskip


Suppose we want to compute the orbits of the projective group 
on the subspaces of the exterior algebra. 
The following command creates $\PGL(5,3)$ 
in the wedge action as in Section~\ref{sec:induced:actions}. 
It then computes the orbits of 
subspaces of the exterior algebra of dimension 5. 
Because of the nature of the poset classification algorithm, 
all orbits of dimension $i$ with $0 \le i \le 5$ are computed. 

\sourcecode{PGL_5_3_wedge_subspace_orbits}

The computation finds 81 orbits at level 5. 
The lex-least spanning tree for the orbits is shown below:

\orbiteroutput{GRAPHICS/WRAPPERS/PGL_5_3_OnWedge_tree_lvl_5}


\bigskip

